
Let $\k$ be a field, and consider the action of the symmetric group $S_n$ on the space $\k[x_1,\dots,x_n]$ of $\k$-valued polynomials by permuting the variables. A polynomial in $\k[x_1,\dots,x_n]$ is \emph{symmetric} if it is invariant under this action. Equivalently, since $S_n$ is generated by transpositions, a~polynomial~$K$ is symmetric if $s_{i_1i_2}K=K$ or $(1-s_{i_1i_2})K=0$ for all $s_{i_1i_2}\in S_n$. One may consider generalizations of symmetric polynomials in which this condition is relaxed, so that we only require $(1-s_{i_1i_2})K$ be divisible by some large polynomial. This leads to the notion of \emph{quasi-invariant polynomials}.

\begin{Definition}
 Let $\k$ be a field. For $m\in\Z_{\geq 0}$, $n\in\Z_{>0}$, a polynomial $K\in\k[x_1,\dots,x_n]$ is \textit{$m$-quasi-invariant} if for all $s_{i_1i_2}\in S_n$ we have that $(x_{i_1}-x_{i_2})^{2m+1}$ divides $(1-s_{i_1i_2})K$. We denote the space of $m$-quasi-invariants by $Q_m(n,\k)$.
\end{Definition}

Note that the symmetric polynomials are exactly the polynomials that are $m$-quasi-invariant for all $m$. For brevity, we also refer to quasi-invariant polynomials as simply quasi-invariants.

Quasi-invariant polynomials were first introduced by Chalykh and Veselov in 1990 \cite{cmp/1104179957} to describe the harmonic, zero eigenvalue eigenfunctions of quantum Calogero--Moser systems. Calogero--Moser systems are a collection of one-dimensional dynamical particle systems that were found to be both solvable \cite{calogero} and integrable \cite{MOSER1975197}. Due to these properties, they have become extensively studied in mathematical physics, with connections to a number of other fields of mathematics, including representation theory.

